% ─────────────────────────────────────────────────────────────────────────────
\section{Introduction}
\label{sec:intro}
% ─────────────────────────────────────────────────────────────────────────────

\subsection{Motivation}

General aviation in the United States operates under a predominantly
\textit{calendar-based} maintenance paradigm: aircraft are inspected at fixed
intervals of hours, days, or calendar cycles regardless of the actual condition
of their components. This reactive approach systematically fails to detect
emerging faults between inspection windows, creating safety risks and generating
costly unscheduled maintenance events that grounding aircraft without warning.

Modern aircraft are equipped with flight data recorders (FDR) and engine control
units (ECU) that continuously log dozens of sensor channels at 1\,Hz or higher.
This operational data, accumulated across thousands of flight hours, encodes the
health history of every critical subsystem. Machine learning offers a principled
path from raw telemetry to \textit{condition-based} maintenance: instead of
waiting for a calendar date, a model can flag anomalous flights and identify the
likely fault category from the sensor profile alone, enabling targeted
interventions before failure occurs.

\subsection{Dataset and Scope}

This project uses the \textbf{National General Aviation Flight Information
Database (NGAFID)} \cite{ngafid2021}, a publicly available benchmark comprising
over 28{,}000 real flights from a Cessna-172 fleet operated by an active flight
school. NGAFID is unique among aviation health datasets in that it provides
(i) authentic multi-sensor telemetry under real operational conditions,
(ii) fault labels derived from verified unscheduled maintenance records
spanning 36 distinct categories, and (iii) a standardized 5-fold cross-validation
split designed to prevent data leakage between flights of the same aircraft
\cite{convtok2023}.

The analysis focuses on the \textit{benchmark 2-day subset}: 11{,}446 flights
across 19 fault categories, preprocessed to fixed length via cubic spline
interpolation to 2{,}048 time steps per flight.

\subsection{Research Questions}

This report investigates two core questions:

\begin{enumerate}
    \item \textbf{Anomaly Detection (AD):} Can a classical ML classifier, trained
          on statistical features extracted from the full flight profile, reliably
          distinguish healthy flights from anomalous ones?
    \item \textbf{Fault Classification (FC):} Given that a flight is anomalous,
          can the same approach identify which of 19 fault types is responsible?
\end{enumerate}

Both tasks are then composed into a single \textbf{End-to-End (E2E) diagnosis
pipeline} that assigns one of 20 labels---\textit{healthy} or one of 19
fault types---to any incoming flight.

\subsection{Hypothesis}

We hypothesize that:

\begin{itemize}
    \item For \textbf{AD}, global statistical features (mean, standard deviation,
          slope, etc.\ over the full flight) are sufficient to characterize the
          overall health state of the aircraft, making classical models competitive
          with deep learning.
    \item For \textbf{FC}, fault signatures manifest as \textit{local micro-patterns}
          in specific temporal segments of the flight profile, contributing less
          than 0.1\% of total variance. Global statistics discard this information,
          creating a structural performance gap versus temporal deep learning models
          that operate directly on the raw time series.
\end{itemize}

\subsection{Contributions}

The main contributions of this work are:

\begin{enumerate}
    \item A systematic evaluation of five classical ML models across multiple
          feature representations (Repr.~B through Repr.~F\textsubscript{++}) on
          the NGAFID benchmark subset, with direct comparison against published
          deep learning results \cite{lmsd2025}.
    \item A feature engineering analysis that identifies cross-physical
          inter-cylinder features (CHT/EGT spreads and ratios) as the most
          informative additions beyond global per-sensor statistics.
    \item A hierarchical fault classification experiment (Oracle K\,=\,4) that
          raises the classical ceiling from F1\,=\,0.121 to F1\,=\,0.420 through
          physics-based fault grouping, closing over half the gap to the DL best
          (MMK~Net, F1\,=\,0.623) without any temporal modeling.
    \item An end-to-end error decomposition showing that 40.6\% of E2E errors
          originate in Stage~2 (wrong fault type), identifying fault classification
          as the bottleneck for future improvement.
\end{enumerate}

\subsection{Report Organization}

The remainder of this report is organized as follows.
Section~\ref{sec:dataset} describes the NGAFID dataset, sensor space, class
distribution, and key data challenges.
Section~\ref{sec:features} presents the feature engineering pipeline and the
hierarchy of representations from global statistics to temporal shape.
Section~\ref{sec:methodology} details the two-stage diagnosis pipeline, the
classical and deep learning models, and the evaluation protocol.
Section~\ref{sec:results} reports and analyzes results for AD, FC, and E2E
across all representations, with comparison to DL benchmarks.
